\documentclass[11pt,a4paper]{article}
% preamble.tex — estilo común de todos los apuntes Froth.
% Cada apunte hace % preamble.tex — estilo común de todos los apuntes Froth.
% Cada apunte hace % preamble.tex — estilo común de todos los apuntes Froth.
% Cada apunte hace \input{preamble} justo después de \documentclass.
\usepackage[margin=2.4cm]{geometry}
\usepackage{xcolor}
\definecolor{litblue}{RGB}{31,79,121}
\definecolor{litgreen}{RGB}{27,94,32}
\definecolor{litorange}{RGB}{230,81,0}
\usepackage{parskip}
\usepackage{titlesec}
\titleformat{\section}{\Large\bfseries\color{litblue}}{\thesection}{0.6em}{}
\titleformat{\subsection}{\large\bfseries\color{litblue}}{\thesubsection}{0.6em}{}
\usepackage{tcolorbox}
\usepackage{fancyhdr}
\pagestyle{fancy}
\fancyhf{}
\fancyhead[L]{\small\color{litblue}Froth — Apuntes de ML}
\fancyhead[R]{\small\thepage}
\renewcommand{\headrulewidth}{0.4pt}
\usepackage[colorlinks=true,linkcolor=litblue,urlcolor=litblue]{hyperref}

% Cabecera de cada apunte: \cabecera{numero}{titulo}
\newcommand{\cabecera}[2]{%
  \begin{tcolorbox}[colback=litblue,coltext=white,colframe=litblue,arc=2mm]
    {\Large\bfseries Apunte #1 — #2}\\[3pt]
    {\small Proyecto Froth · aprender Machine Learning construyendo}
  \end{tcolorbox}\vspace{0.8em}}

% Cajas temáticas
\newtcolorbox{concepto}[1]{colback=blue!4,colframe=litblue,title={Concepto: #1},arc=1.5mm}
\newtcolorbox{practica}{colback=green!5,colframe=litgreen,title={Pruébalo tú (teclado en mano)},arc=1.5mm}
\newtcolorbox{ojo}{colback=orange!8,colframe=litorange,title={Ojo / error típico},arc=1.5mm}
\newtcolorbox{resumen}{colback=gray!8,colframe=gray!60,title={En una frase},arc=1.5mm}

\newcommand{\codigo}[1]{\texttt{#1}}
 justo después de \documentclass.
\usepackage[margin=2.4cm]{geometry}
\usepackage{xcolor}
\definecolor{litblue}{RGB}{31,79,121}
\definecolor{litgreen}{RGB}{27,94,32}
\definecolor{litorange}{RGB}{230,81,0}
\usepackage{parskip}
\usepackage{titlesec}
\titleformat{\section}{\Large\bfseries\color{litblue}}{\thesection}{0.6em}{}
\titleformat{\subsection}{\large\bfseries\color{litblue}}{\thesubsection}{0.6em}{}
\usepackage{tcolorbox}
\usepackage{fancyhdr}
\pagestyle{fancy}
\fancyhf{}
\fancyhead[L]{\small\color{litblue}Froth — Apuntes de ML}
\fancyhead[R]{\small\thepage}
\renewcommand{\headrulewidth}{0.4pt}
\usepackage[colorlinks=true,linkcolor=litblue,urlcolor=litblue]{hyperref}

% Cabecera de cada apunte: \cabecera{numero}{titulo}
\newcommand{\cabecera}[2]{%
  \begin{tcolorbox}[colback=litblue,coltext=white,colframe=litblue,arc=2mm]
    {\Large\bfseries Apunte #1 — #2}\\[3pt]
    {\small Proyecto Froth · aprender Machine Learning construyendo}
  \end{tcolorbox}\vspace{0.8em}}

% Cajas temáticas
\newtcolorbox{concepto}[1]{colback=blue!4,colframe=litblue,title={Concepto: #1},arc=1.5mm}
\newtcolorbox{practica}{colback=green!5,colframe=litgreen,title={Pruébalo tú (teclado en mano)},arc=1.5mm}
\newtcolorbox{ojo}{colback=orange!8,colframe=litorange,title={Ojo / error típico},arc=1.5mm}
\newtcolorbox{resumen}{colback=gray!8,colframe=gray!60,title={En una frase},arc=1.5mm}

\newcommand{\codigo}[1]{\texttt{#1}}
 justo después de \documentclass.
\usepackage[margin=2.4cm]{geometry}
\usepackage{xcolor}
\definecolor{litblue}{RGB}{31,79,121}
\definecolor{litgreen}{RGB}{27,94,32}
\definecolor{litorange}{RGB}{230,81,0}
\usepackage{parskip}
\usepackage{titlesec}
\titleformat{\section}{\Large\bfseries\color{litblue}}{\thesection}{0.6em}{}
\titleformat{\subsection}{\large\bfseries\color{litblue}}{\thesubsection}{0.6em}{}
\usepackage{tcolorbox}
\usepackage{fancyhdr}
\pagestyle{fancy}
\fancyhf{}
\fancyhead[L]{\small\color{litblue}Froth — Apuntes de ML}
\fancyhead[R]{\small\thepage}
\renewcommand{\headrulewidth}{0.4pt}
\usepackage[colorlinks=true,linkcolor=litblue,urlcolor=litblue]{hyperref}

% Cabecera de cada apunte: \cabecera{numero}{titulo}
\newcommand{\cabecera}[2]{%
  \begin{tcolorbox}[colback=litblue,coltext=white,colframe=litblue,arc=2mm]
    {\Large\bfseries Apunte #1 — #2}\\[3pt]
    {\small Proyecto Froth · aprender Machine Learning construyendo}
  \end{tcolorbox}\vspace{0.8em}}

% Cajas temáticas
\newtcolorbox{concepto}[1]{colback=blue!4,colframe=litblue,title={Concepto: #1},arc=1.5mm}
\newtcolorbox{practica}{colback=green!5,colframe=litgreen,title={Pruébalo tú (teclado en mano)},arc=1.5mm}
\newtcolorbox{ojo}{colback=orange!8,colframe=litorange,title={Ojo / error típico},arc=1.5mm}
\newtcolorbox{resumen}{colback=gray!8,colframe=gray!60,title={En una frase},arc=1.5mm}

\newcommand{\codigo}[1]{\texttt{#1}}

\begin{document}
\cabecera{14}{Clustering: que el algoritmo encuentre los grupos (HDBSCAN)}

\section{El problema}

Tu mapa 2D (Apunte 13) muestra nubes a simple vista, pero ``a simple vista'' no es un método:
necesitamos que un algoritmo delimite los grupos de forma objetiva y repetible. Eso es
\textbf{clustering no supervisado}: agrupar sin que nadie diga qué es cada grupo ni cuántos hay.

\section{La intuición clásica: K-means (y sus dos trampas)}

\begin{concepto}{K-means}
El método más famoso. Le dices ``quiero \textbf{K} grupos'' y él reparte los puntos en K
grupos compactos. Simple y rápido, pero con dos trampas para nuestro caso:
\begin{enumerate}
  \item \textbf{Tienes que adivinar K.} ¿Cuántos subtemas tiene tu topic? Justo eso es lo que
        NO sabemos (es lo que queremos descubrir).
  \item \textbf{Obliga a TODOS los puntos a entrar en algún grupo.} El paper de poesía
        acabaría asignado a un cluster de flotación, contaminándolo.
\end{enumerate}
\end{concepto}

\section{La práctica: HDBSCAN, clustering por densidad}

\begin{concepto}{HDBSCAN}
En vez de repartir, HDBSCAN busca \textbf{zonas densas}: regiones donde hay muchos puntos
apretados. Cada zona densa = un cluster. Consecuencias:
\begin{itemize}
  \item \textbf{Decide solo cuántos grupos hay} (nosotros no adivinamos K).
  \item Lo que no pertenece a ninguna zona densa queda como \textbf{ruido} (etiqueta $-1$),
        en vez de contaminar un grupo.
  \item Su parámetro principal es \codigo{min\_cluster\_size}: cuántos puntos mínimo hacen
        falta para que algo cuente como grupo (en Froth: \codigo{MIN\_CLUSTER\_SIZE = 8}).
\end{itemize}
\end{concepto}

\section{Lo que pasó con TUS papers (la historia real)}

Con \codigo{min\_cluster\_size=8}, HDBSCAN encontró \textbf{3 clusters + 9 papers de ruido}
(y con 5, casi idéntico: resultado \emph{estable}, buena señal):

\begin{center}
\begin{tabular}{llp{6.6cm}}
\textbf{Grupo} & \textbf{Tamaño} & \textbf{Qué resultó ser (Apunte 15)} \\
\hline
Cluster 0 & 49 & Geología de pegmatitas y granitos de metales raros \\
Cluster 2 & 43 & Flotación de lepidolita (el corazón de la tesis) \\
Cluster 1 & 25 & ``Cajón de lo ajeno'': reciclaje/economía del litio + los off-topic \\
Ruido & 9 & Puentes generalistas y rarezas \\
\end{tabular}
\end{center}

\begin{ojo}
\textbf{Dos sorpresas que enseñan mucho:}
\begin{enumerate}
  \item \textbf{``Ruido'' no siempre es basura.} En el ruido cayeron dos clásicos totalmente
        on-topic (p.ej.\ \emph{The Effect of Bubble Size on Fine Particle Flotation}, 527
        citas). ¿Por qué? Son papers \emph{generalistas} que hablan con varios subtemas a la
        vez: viven ``entre'' las nubes, sin pertenecer del todo a ninguna.
  \item \textbf{Los off-topic no fueron ruido: fueron puestos en cuarentena juntos.} El paper
        de telecomunicaciones, el de poesía y el de historia cayeron los tres dentro del
        cluster 1. ``No ser del tema'' también es algo en común, y forma su propia nube lejos
        de las nubes mineras.
\end{enumerate}
\end{ojo}

\begin{resumen}
K-means exige adivinar K y fuerza a todos adentro; HDBSCAN agrupa por densidad, decide solo
cuántos grupos hay y aparta el ruido. En tus datos: 3 subtemas estables + 9 ruidos, con los
off-topic en cuarentena dentro de su propio cluster.
\end{resumen}

\end{document}
